% M2 proof module: audited WP11/WP13/WP19 arbitrary-support autonomous bridge.

\subsection{Arbitrary coherent-support mixed survival--synthesis bridge}
\label{supp:mixed-full-proof}

Let $A_\nu$ satisfy the exact exchange relations
\begin{equation}
[H_S,A_\nu]=+\hbar\nu A_\nu,
\qquad
[H_C,A_\nu]=-\hbar\nu A_\nu,
\end{equation}
and let $P=\supp\rho_0$, $Q=\id-P$. A two-sided differentiable physical family implies
\begin{equation}
QA_\nu Q=0.
\end{equation}
Decompose
\begin{equation}
B=PA_\nu P,
\qquad
K_+=QA_\nu P,
\qquad
K_-=QA_\nu^\dagger P.
\label{eq:mixed-proof-decomp}
\end{equation}
Then
\begin{equation}
A_\nu=B+K_++K_-^\dagger.
\end{equation}
Define
\begin{align}
J_B^+&=\Tr(B\rho_0^+B^\dagger),&
J_B^-&=\Tr(B^\dagger\rho_0^+B),\\
J_+&=\Tr(K_+\rho_0^+K_+^\dagger),&
J_-&=\Tr(K_-\rho_0^+K_-^\dagger).
\end{align}

\subsubsection{Measurement-side bound}

Set
\begin{equation}
X=(B+K_+)P,
\qquad
Y=K_-.
\end{equation}
The outputs of $B$ and $K_+$ lie in orthogonal $P$ and $Q$ sectors, hence
\begin{equation}
\Tr(X\rho_0^+X^\dagger)=J_B^++J_+.
\end{equation}
For a POVM, the complex score is the sum of the right-supported $X$ score and the conjugate $Y$ score. Weighted Hilbert--Schmidt Cauchy--Schwarz and Minkowski therefore give
\begin{equation}
\sqrt{\Tr F_1}
\le
\sqrt{J_B^++J_+}+\sqrt{J_-}.
\end{equation}
Applying the same argument to $A_\nu^\dagger$ gives
\begin{equation}
\sqrt{\Tr F_1}
\le
\sqrt{J_B^-+J_-}+\sqrt{J_+}.
\end{equation}
The independent-copy weighted norms scale by $N$, so for every finite $N$ and arbitrary collective POVM,
\begin{equation}
\boxed{
\sqrt{\frac{\Tr F_N}{N}}
\le
\min\!\left\{
\sqrt{J_B^++J_+}+\sqrt{J_-},
\sqrt{J_B^-+J_-}+\sqrt{J_+}
\right\}.
}
\label{eq:mixed-proof-measurement}
\end{equation}

\subsubsection{Two-sided shorted survival ceilings}

Because $P$ need not commute with either local Hamiltonian, $B=PAP$ is generally not itself an exact Bohr-gap operator. The energy resource of $B$ therefore depends on the angle between its information-bearing support and the physical endpoint subspaces.

Let
\begin{equation}
Z_B^+=B\rho_0^+B^\dagger,
\qquad
Z_B^-=B^\dagger\rho_0^+B,
\end{equation}
and
\begin{equation}
R_B^+=\supp Z_B^+,
\qquad
R_B^-=\supp Z_B^-.
\end{equation}
Let $\Pi_{S,U},\Pi_{S,D},\Pi_{C,U},\Pi_{C,D}$ be the minimal local endpoint projectors containing the participating support-preserving transition content. Define compressed endpoint operators
\begin{align}
S_{S,U}&=P\Pi_{S,U}P,&S_{S,D}&=P\Pi_{S,D}P,\\
S_{C,U}&=P\Pi_{C,U}P,&S_{C,D}&=P\Pi_{C,D}P.
\end{align}
The four shorting constants are
\begin{align}
\lambda_{S,U}&=\sup\{\lambda:S_{S,U}\succeq\lambda R_B^+\},\\
\lambda_{C,D}&=\sup\{\lambda:S_{C,D}\succeq\lambda R_B^+\},\\
\lambda_{S,D}&=\sup\{\lambda:S_{S,D}\succeq\lambda R_B^-\},\\
\lambda_{C,U}&=\sup\{\lambda:S_{C,U}\succeq\lambda R_B^-\}.
\end{align}
Let $R_B$ be the physical affine radius of the support-preserving sub-tangent $B$ and $T_{X,E}=\Tr(\Pi_{X,E}\rho_0)$. The finite-radius operator majorization gives
\begin{align}
J_B^+&\le\frac{4T_{S,U}}{R_B^2\lambda_{S,U}},&
J_B^+&\le\frac{4T_{C,D}}{R_B^2\lambda_{C,D}},\\
J_B^-&\le\frac{4T_{S,D}}{R_B^2\lambda_{S,D}},&
J_B^-&\le\frac{4T_{C,U}}{R_B^2\lambda_{C,U}}.
\end{align}
Thus define
\begin{align}
a_+&=
\min\!\left\{
\frac{4T_{S,U}}{R_B^2\lambda_{S,U}},
\frac{4T_{C,D}}{R_B^2\lambda_{C,D}}
\right\},\\
a_-&=
\min\!\left\{
\frac{4T_{S,D}}{R_B^2\lambda_{S,D}},
\frac{4T_{C,U}}{R_B^2\lambda_{C,U}}
\right\},
\end{align}
so
\begin{equation}
J_B^+\le a_+,
\qquad
J_B^-\le a_-.
\label{eq:mixed-proof-apm}
\end{equation}
If $B=0$, set $a_+=a_-=0$.

\subsubsection{Canonical endpoint-incidence action}

Define the canonical joint endpoint-role projectors of the full exact exchange,
\begin{equation}
\Pi_{\rm out}=\supp(A_\nu A_\nu^\dagger),
\qquad
\Pi_{\rm in}=\supp(A_\nu^\dagger A_\nu).
\end{equation}
The exchange commutators imply
\begin{equation}
[H_X,A_\nu A_\nu^\dagger]=0,
\qquad
[H_X,A_\nu^\dagger A_\nu]=0
\end{equation}
for $X=C,S$, hence the two support projectors commute with both local Hamiltonians. They encode the output and input endpoint roles of the full ladder.

One joint endpoint incidence corresponds to one absolute signal gap and one absolute clock gap. Define
\begin{equation}
G_{\rm ex}
=2\hbar\nu\,Q(\Pi_{\rm out}+\Pi_{\rm in})Q.
\label{eq:mixed-proof-Gex}
\end{equation}
If a joint energy state is both a domain and range endpoint, it is counted twice because it plays two distinct roles.

Let
\begin{equation}
C_\Delta=Q(\partial_x^2+\partial_y^2)\rho(0)Q.
\end{equation}
Second-order positivity gives
\begin{equation}
C_\Delta\succeq Z_++Z_-,
\end{equation}
where
\begin{equation}
Z_+=K_+\rho_0^+K_+^\dagger,
\qquad
Z_-=K_-\rho_0^+K_-^\dagger.
\end{equation}
Define
\begin{equation}
\Aact_{\rm ex}^{(2)}
=\frac14\Tr(G_{\rm ex}C_\Delta).
\end{equation}
Let $R_\pm=\supp Z_\pm$ and
\begin{equation}
g_\pm=
\lambda_{\min}
(R_\pm G_{\rm ex}R_\pm|_{R_\pm})
\end{equation}
for nonzero components. Then
\begin{align}
4\Aact_{\rm ex}^{(2)}
&=\Tr(G_{\rm ex}C_\Delta)\\
&\ge\Tr(G_{\rm ex}Z_+)+\Tr(G_{\rm ex}Z_-)\\
&\ge g_+J_+ +g_-J_-.
\label{eq:mixed-proof-action-budget}
\end{align}
The same physical curvature is charged once.

\subsubsection{Exact scalar reduction}

For $a,e\ge0$ and $p,q>0$, maximize
\begin{equation}
\left[\sqrt{a+j_+}+\sqrt{j_-}\right]^2
\end{equation}
subject to
\begin{equation}
pj_+ +qj_-\le e,
\qquad j_+,j_-\ge0.
\end{equation}
Because the objective is increasing, the budget is saturated at the optimum. Set
\begin{equation}
j_-=\frac{e-pj_+}{q},
\qquad0\le j_+\le e/p.
\end{equation}
The derivative of the square-root sum vanishes when
\begin{equation}
\frac1{\sqrt{a+j_+}}
=\frac{p/q}{\sqrt{(e-pj_+)/q}},
\end{equation}
which yields an interior solution exactly when
\begin{equation}
e\ge\frac{ap^2}{q}.
\end{equation}
Evaluating the boundary and interior branches gives
\begin{equation}
\Psi_a(e;p,q)=
\begin{cases}
(\sqrt a+\sqrt{e/q})^2,&e\le ap^2/q,\\[1mm]
(e+pa)(p^{-1}+q^{-1}),&e\ge ap^2/q.
\end{cases}
\label{eq:mixed-proof-Psi}
\end{equation}

Use Eq.~\eqref{eq:mixed-proof-apm} and Eq.~\eqref{eq:mixed-proof-action-budget} in the first orientation of Eq.~\eqref{eq:mixed-proof-measurement} to obtain
\begin{equation}
\frac{\Tr F_N}{N}
\le
\Psi_{a_+}(4\Aact_{\rm ex}^{(2)};g_+,g_-).
\end{equation}
The conjugate orientation gives
\begin{equation}
\frac{\Tr F_N}{N}
\le
\Psi_{a_-}(4\Aact_{\rm ex}^{(2)};g_-,g_+).
\end{equation}
Therefore
\begin{equation}
\boxed{
\frac{\Tr F_N}{N}
\le
\min\!\left\{
\Psi_{a_+}(4\Aact_{\rm ex}^{(2)};g_+,g_-),
\Psi_{a_-}(4\Aact_{\rm ex}^{(2)};g_-,g_+)
\right\}.
}
\label{eq:mixed-proof-final}
\end{equation}

When $K_+=K_-=0$, the action term vanishes and Eq.~\eqref{eq:mixed-proof-final} reduces to two-sided finite-radius survival. When $B=0$ and the endpoint geometry is clean, $g_+=g_-=2\hbar\nu$ and the bilateral harmonic price is $\hbar\nu$, recovering the sharp autonomous boundary coefficient $\hbar\nu/4$.

\subsubsection{Shared-kernel qutrit check}

For the fixed-total-excitation-2 qutrit benchmark,
\begin{equation}
A=|M\rangle\langle L|-\sqrt2|U\rangle\langle M|,
\end{equation}
so
\begin{equation}
\Pi_{\rm out}=|M\rangle\langle M|+|U\rangle\langle U|,
\end{equation}
\begin{equation}
\Pi_{\rm in}=|L\rangle\langle L|+|M\rangle\langle M|.
\end{equation}
Thus
\begin{equation}
G_{\rm ex}/(\hbar\nu)=\operatorname{diag}(2,4,2).
\end{equation}
The exact weighted norms are
\begin{equation}
J_B^+=J_B^-=\frac54,
\qquad
J_+=\frac74,
\qquad
J_-=3.
\end{equation}
Both synthesized ranges are the same rank-one kernel direction, giving
\begin{equation}
g_+=g_-=\frac{13}{4}\hbar\nu.
\end{equation}
For the minimal curvature,
\begin{equation}
4\Aact_{\rm ex}^{(2)}=\frac{247}{16}\hbar\nu.
\end{equation}
Using $a_+=a_-=5/4$ and units $\hbar\nu=1$, the second branch of Eq.~\eqref{eq:mixed-proof-Psi} gives
\begin{equation}
\Psi_{5/4}(247/16;13/4,13/4)=12.
\end{equation}
Hence the canonical autonomous action reduction introduces no additional resource-layer looseness in this benchmark.
